\section{总结与延伸}

\subsection{本章小结}

这门课最后一讲想留下的，不只是几个模型名字，而是一套判断框架。

\begin{importantbox}{Lecture 16 的核心 takeaway}
\begin{enumerate}
    \item foundation model 的本质是大规模、通用、可迁移的预训练表示。
    \item CLIP 用图文对比学习把图像和文本放进统一空间。
    \item zero-shot 分类可以直接把文本当成类别原型。
    \item scale 对 CLIP 的成功至关重要，尤其是数据和参数。
    \item CoCa 说明在对比学习之外，加上生成目标也很有价值。
    \item LLaVA 和 Flamingo 把视觉 token 接进 LLM，打开了多模态对话和 few-shot reasoning。
    \item grounding、retrieval 和 verifier 是降低幻觉、提升系统可靠性的关键手段。
    \item 真正强的系统往往是 chaining，而不是单模型裸奔。
\end{enumerate}
\end{importantbox}

\begin{knowledgebox}{这一讲最值得记住的系统观}
模型越强，系统设计越不能简化。视觉 foundation model 的价值，不只是“让模型会更多”，而是“让模型更容易被组合、验证和复用”。
\end{knowledgebox}

\begin{warningbox}{最现实的限制}
开放模型、开放数据、开放代码虽然重要，但真正“可复现”的能力仍然很难。很多最强的视觉语言模型，外界只能看到结果，看不到完整训练配方。这个缺口本身就是未来研究的主战场。
\end{warningbox}

\subsection{延伸阅读}

\begin{itemize}
    \item Radford et al., \emph{Learning Transferable Visual Models From Natural Language Supervision}
    \item Yu et al., \emph{CoCa: Contrastive Captioners are Image-Text Foundation Models}
    \item Liu et al., \emph{LLaVA}
    \item Alayrac et al., \emph{Flamingo}
    \item Winoground / ARO / CREPE 等 compositionality benchmarks
    \item CuPL: prompting CLIP-style models with richer class descriptions
\end{itemize}

\end{document}
